% !TeX root = ../教材题库_全部条目_转写工作稿.tex
% Source: XJ-B1-C05-TH-P177-Y0904
% 第5章 / 5.3 三角函数的图象与性质 / 教材 p.177、178 / PDF 第184、185页
% OCR未校对：公式与文字均可能有误；未完成数学TeX转写；禁止直接发给学生
% 原题及相邻原解截图（核对用）：../题目原图/XJ-B1-C05-TH-P177-Y0904-01.webp, ../题目原图/XJ-B1-C05-TH-P177-Y0904-02.webp
% 本文件只在转写完成后作为题干使用；OCR 条目尚不是可靠题干。
\begin{bankitem}{p.177 思考（含紧邻说明）}{第5章 / 5.3 三角函数的图象与性质 / 教材 p.177、178 / PDF 第184、185页}{XJ-B1-C05-TH-P177-Y0904}{OCR未校对：公式与文字均可能有误；未完成数学TeX转写；禁止直接发给学生}
\begin{ocrdraft}
[教材p.177 / PDF第184页]
\par
？思考如何利用正弦线画出y=sinx在区间[0，2冠上的图象？
\par
如图5．3一1，在平面直角坐标系的轴上取一点0，以0为圆心画单位圆交工
\par
轴于点A，且有向线段OIA指向轴的正方向．在单位圆上任取一点P,作PD垂
\par
直于轴，垂足为D,则DP是对应角ZPOIA的正弦线，DP=sin“，其中“为
\par
ZPOIA的弧度数．
\par
M(xo,sin）
\par
0
\par
0
\par
图5．3一1
\par
在轴上取点N()o，0），将正弦线DP平移至NM，使点D与N重合，可知
\par
点M的坐标为()o，sin“），它是函数y=sinx图象上的一点．
\par
随着弧度数的改变，可以得到函数y=sin图象上的其他点．为方便起见，
\par
可将圆0平均分成12等份（分得越多，相应作出的图象越精确），使“的值依次取
\par
\ensuremath{\cdot}\ensuremath{\cdot}，2丌，借助于它们的正弦线，依次作出函数3-艹图象上的点
\par
[教材p.178 / PDF第185页]
\par
（0，sin0），
\par
，S1n
\par
，S1n
\par
，S1n
\par
（2丌，sin2丌），用一条光
\par
滑的曲线将这些点依次连接起来，就得到函数3，一sin在区间@，2丌]上的图象，如
\par
图5．3一2．
\par
7芤43丌5丌ll丌
\par
是
\par
丌
\par
是2芤5是
\par
2兀
\par
），=sin工，xe[0，2]
\par
图5．3一2
\par
由于sin(x十2丌）一艹对所有整数都成立，因此，将3，=sinx在区间@，
\par
2丌]上的图象逐次向左和向右平移2丌个单位长度，就可以得到正弦函数3，一sin产
\par
xeR的图象（如图5，3一3）．
\par
y=Sln.X
\par
图5．3一3
\par
正弦函数3，一sin的图象称为正弦曲线．
\par
我们可以借助计算机作图软件快捷精准地画出三角函数在指定区间上的图象，
\par
由于正弦曲线在区间[0，2丌]上有五个点（最高点、最低点、与工轴的交点）（0，0），
\par
3丌
\par
，（2丌，0）对曲线的升降起伏起着关键作用，于是在精
\par
度要求不太高的时候，只要画出了这五个点，曲线的大致形状就基本确定了，将它们依
\par
次连成光滑曲线，就得到正弦曲线的简图．正弦曲线的这种近似画法称为“五点法”．
\end{ocrdraft}
\end{bankitem}
